\documentclass[11pt]{article}
\usepackage[utf8]{inputenc}
\usepackage[T1]{fontenc}
\usepackage{lmodern}
\usepackage{amsmath,amssymb,amsthm}
\usepackage[margin=1in]{geometry}
\usepackage{booktabs,tabularx}
\usepackage[hidelinks]{hyperref}
\usepackage{url}
\usepackage{array}
\def\UrlBreaks{\do0\do1\do2\do3\do4\do5\do6\do7\do8\do9\do a\do b\do c\do d\do e\do f\do/\do.\do-}
\emergencystretch=3em
\newtheorem{theorem}{Theorem}
\newtheorem{lemma}[theorem]{Lemma}
\newtheorem{corollary}[theorem]{Corollary}
\newtheorem{conjecture}[theorem]{Conjecture}
\theoremstyle{remark}\newtheorem{remark}[theorem]{Remark}
\newcommand{\lean}[1]{{\ttfamily\def\_{\textunderscore\allowbreak}\def\/{/\allowbreak}#1}}
\author{Hetark Shah, AI-assisted by Grok}
\date{October 12, 2026}
\title{On Ordowski's characterization of odd primes by the number of squares modulo $n$ (OEIS A000224): even numbers, prime powers and semiprimes, formalized in Lean~4}
\begin{document}
\maketitle
\begin{abstract}
Let $a(n)$ be the number of squares modulo $n$ (OEIS A000224~\cite{oeis}). T.~Ordowski conjectured that for $n>1$, $n^2\equiv1\pmod{a(n)(a(n)-1)}$ holds if and only if $n$ is an odd prime. \textbf{The full conjecture remains open.} We prove it for every $n>1$ that is even, a prime power $p^k$ ($k\ge1$), or a product $pq$ of two distinct odd primes. The prime-power case uses Stangl's formula for $a(p^k)$~\cite{stangl}, which we also formalize. The semiprime case combines a size bound with a divisibility bound and a finite check. Everything is formalized in Lean~4 with Mathlib against the exact definition of $a$ in Google DeepMind's \texttt{formal-conjectures} repository~\cite{fc}. The development uses no \texttt{sorry} and no \texttt{native\_decide}, and only the standard axioms.
\end{abstract}

\section{Introduction}
For $n\ge1$ let $a(n)=\#\{k^2\bmod n:0\le k<n\}$.
\begin{conjecture}[Ordowski, 2025]\label{conj}
For $n>1$, $n$ is an odd prime if and only if $n^2\equiv1\pmod{a(n)(a(n)-1)}$.
\end{conjecture}
The conjecture appears as a comment on A000224~\cite{oeis}, together with M.~Marcus's remark that it holds up to $n=10^8$. It is formalized as \lean{OeisA224.conjecture} in \nolinkurl{FormalConjectures/OEIS/224.lean} of \texttt{formal-conjectures}~\cite{fc}, where it is tagged \texttt{research open}. The forward direction is immediate from $a(p)=(p+1)/2$. The open part is to show that no composite $n$ satisfies the congruence.

\section{Related work}
The count $a(n)$ is multiplicative, and Stangl~\cite{stangl} gives its value at prime powers; see also Finch and Sebah~\cite{fs}. On October 12, 2026 we searched OEIS, the web, MathOverflow and the issues and pull requests of \texttt{formal-conjectures} (upstream commit \texttt{34ee0e3}). We found no proof of Conjecture~\ref{conj}, partial or complete. The cases treated here therefore appear to be new, although the arguments are elementary. The formalization uses Lean~4~\cite{lean4} and Mathlib~\cite{mathlib}.

\section{Statement}
\begin{theorem}\label{main}
Conjecture~\ref{conj} holds for every $n>1$ that is (i) even, (ii) a prime power $p^k$ with $k\ge1$, or (iii) a product $pq$ of distinct odd primes $p,q$.
\end{theorem}
\noindent Cases not covered include squarefree $n$ with three or more odd prime factors, and odd $n$ such as $p^2q$ that have a repeated prime factor and at least two distinct primes.

\section{Proof overview}
\subsection{Preliminaries}
\begin{lemma}\label{mult}
$a(n)=\#\{x^2 : x\in\mathbb{Z}/n\}$; $a$ is multiplicative; and $a(p)=(p+1)/2$ for odd primes $p$.
\end{lemma}
\begin{proof}
The first claim holds because reduction mod $n$ is a bijection $\{0,\dots,n-1\}\to\mathbb{Z}/n$. If $\gcd(m,n)=1$, the CRT ring isomorphism $\mathbb{Z}/mn\cong \mathbb{Z}/m\times\mathbb{Z}/n$ commutes with squaring, so it maps the set of squares onto the product of the two sets of squares. For an odd prime $p$, the squares are the values $k^2$ with $0\le k\le (p-1)/2$, since $x^2=(-x)^2$. These values are distinct: $k^2=l^2$ forces $p\mid (k-l)(k+l)$, and $0\le k+l<p$.
\end{proof}

\subsection{The even case}
If $n$ is even, then $n^2-1$ is odd, while $a(n)(a(n)-1)$ is a product of two consecutive integers and hence even. So $n^2\not\equiv 1$. (If $a(n)(a(n)-1)=0$, the congruence would say $n^2=1$, which is impossible for $n>1$.)

\subsection{The semiprime case}
\begin{theorem}\label{pq}
Let $p<q$ be odd integers with $p\ge 3$ prime, and let $s=(p+1)(q+1)/4$. Then $s(s-1)\nmid (pq)^2-1$.
\end{theorem}
Primality of $q$ is not needed. By Lemma~\ref{mult}, $a(pq)=s$ for distinct odd primes, so Theorem~\ref{pq} gives the semiprime case of Theorem~\ref{main}.
\begin{proof}
Put $P=(p+1)(q+1)$, so that $16\,s(s-1)=P(P-4)$. Suppose $Q\,s(s-1)=(pq)^2-1$, so $Q\,P(P-4)=16((pq)^2-1)$.

\emph{Step 1: $Q\le 15$.} Expanding, $P(P-4)-((pq)^2-1)=2pq(p+q+1)-4pq+(p+q+1)^2-4(p+q+1)+1>0$. Hence $Q<16$.

\emph{Step 2: $p\le 61$.} By Step~1, $16((pq)^2-1)\le 15P(P-4)<15P^2$. Suppose $p\ge 62$. Since $q>p$ we have $(q+1)p\le (p+1)q$, so $P^2p^2\le (p+1)^4q^2$. Also $15(p+1)^4\le 16p^4-16$ for $p\ge 62$. Together with $q\ge p$ this gives $15P^2\le 16((pq)^2-1)$, a contradiction.

\emph{Step 3: $q+1\mid 16(p^2-1)$.} Since $q+1\mid P$, we get $q+1\mid 16((pq)^2-1)$. Also $16((pq)^2-1)-16(p^2-1)=16p^2(q^2-1)$ is divisible by $q+1$.

\emph{Step 4: finite check.} For each prime $p\in[3,61]$ and each divisor $d=q+1$ of $16(p^2-1)$ with $q>p$ odd, direct computation shows $s(s-1)\nmid (pq)^2-1$. There are only finitely many such pairs.
\end{proof}

\subsection{Odd prime powers}
\begin{lemma}[Stangl~\cite{stangl}]\label{stangl}
For an odd prime $p$ and $k\ge 2$, $a(p^k)=a(p^{k-2})+\tfrac12 p^{k-1}(p-1)$. Consequently $2(p+1)\,a(p^k)=p^{k+1}+p+2$ if $k$ is even and $p^{k+1}+2p+1$ if $k$ is odd.
\end{lemma}
\begin{proof}
Let $N=p^k$. A square $x^2 \bmod N$ is divisible by $p$ iff $p\mid x$. Writing $x=pt$ gives $(pt)^2 \bmod p^2M=p^2(t^2\bmod M)$ with $M=p^{k-2}$, so the squares divisible by $p$ are exactly $p^2$ times the squares modulo $M$. For the squares prime to $p$: if $x,y\in[0,N/2)$ are prime to $p$ and $N\mid (y-x)(y+x)$, then $p$ cannot divide both factors (as $p$ is odd and $p\nmid y$), so $N$ divides one factor, forcing $x=y$. Replacing $x$ by $N-x$, every unit square has such a representative. Hence the unit squares number $\varphi(N)/2$. The closed form follows by induction from $a(1)=1$ and $a(p)=(p+1)/2$.
\end{proof}
\begin{theorem}\label{pp}
If $p$ is an odd prime, $k\ge 2$, and $s$ satisfies the identity of Lemma~\ref{stangl}, then $s(s-1)\nmid p^{2k}-1$. Consequently, $(p^k)^2\not\equiv 1\pmod{a(p^k)(a(p^k)-1)}$, and Conjecture~\ref{conj} holds for all odd prime powers.
\end{theorem}
\begin{proof}[Sketch]
Let $X=p^k$ and $M=2(p+1)$, and write $Q\,s(s-1)=X^2-1$. Multiplying by $M^2$ gives $Q(Ms)(Ms-M)=M^2(X^2-1)$. For even $k$, reducing modulo $p$ gives $p\mid 4(p+1)^2$, which is impossible. For odd $k\ge 3$, a size estimate gives $1\le Q\le 7$, and reducing modulo $p$ gives $p\mid Q-4$. If $Q=4$, the equation becomes $(2p+1)X^2=2p^2X+p^2$, which is impossible since $X>p$. Otherwise $p=3$ and $Q\in\{1,7\}$, and both cases are ruled out by elementary inequalities.
\end{proof}



\section{Formalization}
All declarations are in namespace \lean{OeisA224}. The definition \lean{a} in \lean{Defs.lean} is copied verbatim from \texttt{formal-conjectures}. The theorem \lean{conjecture\_main} concludes exactly \lean{(n.Prime \(\wedge\) n \(\neq\) 2) \(\leftrightarrow\) n \textasciicircum{} 2 \(\equiv\) 1 [MOD a n * (a n - 1)]} under its case hypothesis.
\begin{table}[ht]
\centering
\begin{tabularx}{\linewidth}{@{}>{\raggedright\arraybackslash}X>{\raggedright\arraybackslash}X@{}}
\toprule
Lean name (file) & Statement \\
\midrule
\lean{a\_eq\_aZ}, \lean{aZ\_mul}, \lean{aZ\_prime} (\lean{Count.lean}) & Lemma~\ref{mult} \\
\lean{even\_case}, \lean{prime\_direction} (\lean{Basic.lean}) & even case; odd prime $\Rightarrow$ congruence \\
\lean{finCheck\_true} (\lean{FiniteCheck.lean}) & Step~4, checked with \lean{decide +kernel} \\
\lean{pq\_reduction}, \lean{pq\_case}, \lean{semiprime\_case} (\lean{Semiprime.lean}) & Theorem~\ref{pq} and case (iii) \\
\lean{pp\_even}, \lean{pp\_odd} (\lean{PrimePowerCore.lean}) & arithmetic core of Theorem~\ref{pp} \\
\lean{stangl}, \lean{prime\_power\_case} (\lean{PrimePower.lean}) & Lemma~\ref{stangl}; Theorem~\ref{pp} \\
\lean{conjecture\_of\_even\_or\_prime\_or\_semiprime}, \lean{conjecture\_main} (\lean{Main.lean}) & Theorem~\ref{main} \\
\bottomrule
\end{tabularx}
\caption{Theorems and their Lean names.}\label{tab:lean}
\end{table}

\section{Reproducibility}
\begin{itemize}
\item Source code: \url{https://github.com/het4rk/ordowski-a224}.
\item Toolchain: \texttt{leanprover/lean4:v4.35.0-rc4}. Mathlib: tag \texttt{v4.35.0-rc4}, commit \nolinkurl{1f414401f69059aa7eead47b53ee40bd38455eeb}, pinned in \texttt{lake-manifest.json}.
\item Build:
\begin{verbatim}
lake exe cache get
lake build
lake env lean Axioms.lean
\end{verbatim}
 \lean{FiniteCheck.lean} takes a few minutes and several GB of RAM.
\item Axioms: all 12 declarations listed in \texttt{Axioms.lean}, including \lean{conjecture\_main}, depend only on \texttt{propext}, \texttt{Classical.choice} and \texttt{Quot.sound}.
\item Attestation: \texttt{verification/} contains the fresh-clone build log, the axiom output and \texttt{VERIFICATION.md}, which records the commit, the toolchain and SHA-256 hashes of the Lean sources. CI (\url{https://github.com/het4rk/ordowski-a224/actions/workflows/lean.yml}) rebuilds the project on every push. It rejects \texttt{sorry} and non-whitelisted axioms, and uploads the build log and axiom report as artifacts.
\end{itemize}

\section{Limitations}
\textbf{The full conjecture remains open.} We do not treat squarefree $n$ with three or more prime factors, or general odd composites with a repeated factor and two or more distinct primes (for example $p^2q$). One part of the method extends but is not formalized: the argument of Step~3 shows that if $n=mq$ satisfies the congruence, with $q$ prime and $q\nmid m$, then $q+1\mid4(m^2-1)$. The size argument of Step~2, however, does not bound the second-smallest prime once there are three or more factors, so the method does not directly extend. The independent computational checks in \texttt{scripts/} are not part of the formal proof.

\begin{thebibliography}{9}
\bibitem{oeis} OEIS Foundation Inc., \emph{A000224: Number of squares mod n}, \url{https://oeis.org/A000224}.
\bibitem{fc} Google DeepMind, \emph{formal-conjectures}, \url{https://github.com/google-deepmind/formal-conjectures}, file \nolinkurl{FormalConjectures/OEIS/224.lean}.
\bibitem{stangl} W.~D.~Stangl, Counting squares in $\mathbb{Z}_n$, \emph{Math.\ Mag.} \textbf{69} (1996), 285--289.
\bibitem{fs} S.~R.~Finch, P.~Sebah, \emph{Squares and cubes modulo $n$}, arXiv:math/0604465.
\bibitem{mathlib} The mathlib Community, \emph{The Lean mathematical library}, in Proc.\ CPP 2020, ACM, 367--381.
\bibitem{lean4} L.~de Moura, S.~Ullrich, \emph{The Lean 4 theorem prover and programming language}, in CADE-28, LNCS 12699, Springer, 2021, 625--635.
\end{thebibliography}
\end{document}
